\section{Residual-matched rational terms on the original window}
\label{sec:rational-one-prime-actual-window}

\paragraph{The exact preceding remainder.}
Keep every original finite datum, page, prescribed pair, coefficient
completion and strict local condition of AC-G192-window.
Thus $h\geq2$, $T\geq1$ are integers and, at source size $u$,
\[
 \begin{gathered}
 H_u=\mathbb Q(\zeta_{2^{h+2}},\zeta_7,\sqrt{-q},\sqrt{-d},
                         \omega,\sqrt[3]{269}),\quad K=\mathbb Q(\omega),\\
 u^{2+8\varepsilon}\leq2^h\leq2u^{2+8\varepsilon},\qquad
 T=\lfloor u^{399}\rfloor,\qquad0<\varepsilon<1/256.
 \end{gathered}\tag{RMW.1}
\]
The distinct auxiliary primes satisfy $q,d>7$ and
$q,d\equiv3\pmod4$, exactly as in the retained source.
Let $Z_u$, $J_u=J_{u,0}\sqcup J_{u,\rm mix}$ and $V_u$
be precisely the sets of (CAW.1)--(CAW.4).
In particular $Z_u$ is the complete G190 omitted union, and
$V_u$ consists only of the two original-residue conductor-$32$
coefficient-three rows specified in (CAW.3).
For a trivial-cubic label, use the unique rational finite extension
$\psi_{\eta,e}^{\rm nc}$ with sign $(-1)^e$, on the same original
$O_\ell$. The rational strict group is full at $\ell$, zero on
other finite inertia, and has the real condition of AC-G191.

\paragraph{The residual-matched actual labels.}
Let $\tau_7$ be the Teichm\"uller lift of $\overline c_7$ in
$\mathbb Z_7^\times$, and let $\lambda_r$ be the primitive
Legendre character of conductor $r$, for $r=q,d$.
Define, with every undisplayed finite factor trivial,
\[
 \begin{aligned}
 \mathfrak P_+(T)&=\{\ell\leq T:\ell\geq7\text{ prime},\
                    \ell\equiv1\pmod3,\ \ell\ne7\},\\
 Y_+(u)&=\{(\ell,1,e):\ell\in\mathfrak P_+(T),\
                       1\leq e<T,\ (\ell-1)\mid e\},\\
 Y_7(u)&=\{(7,\tau_7^e|_{G_K},e):1\leq e<T\},\\
 Y_r(u)&=\{(r,\lambda_r|_{G_K},e):1\leq e<T,\
                   e\equiv(r-1)/2\pmod{r-1}\},\quad r=q,d,\\
 Y_u&=Y_+(u)\sqcup Y_7(u)\sqcup Y_q(u)\sqcup Y_d(u).
 \end{aligned}\tag{RMW.2}
\]
These are the literal families of (RQM.4), with the original
coefficient embeddings and without any extra continuous twist.
Set
\[
 \begin{split}
 \mathcal R_{399}^{\rm match}(u)={}&
 \sum_{(\ell,\eta,e)\in J_{u,\rm mix}}
                  \frac{\log\#\mathcal S_{\ell,\eta}(e)}{b_\ell}\\
 &+\sum_{(\ell,\eta,e)\in J_{u,0}\setminus(V_u\cup Y_u)}
       \frac{\log\#\mathcal S^{\mathbb Q}_{\ell,
                              \psi_{\eta,e}^{\rm nc}}(e)}{b_\ell}.
 \end{split}\tag{RMW.3}
\]
No mixed cubic term is removed here. No original positive index
is identified with a sum of rational positive indices.

\begin{theorem}[AC-G196-window: residual matching on the same window]
\label{AC-G196-window}
The sets in (RMW.2) are internally disjoint actual triples, and
\[
 Y_u\subset J_{u,0}\setminus V_u,\qquad
 \mathcal R_{399}^{\rm match}(u)=\mathcal R_{399}^{\rm nc}(u).
 \tag{RMW.4}
\]
The equality is exact: every deleted rational noncritical group
is zero on its original coefficient. Its critical partner is
retained in the already proved single critical-cost aggregate.
For $N_T=T-1$, put
\[
 \begin{aligned}
 J_+(T)&=\sum_{\ell\in\mathfrak P_+(T)}
                         \left\lfloor\frac{T-1}{\ell-1}\right\rfloor,\\
 A_r(T)&=\left\lfloor\frac{T-1+(r-1)/2}{r-1}\right\rfloor
                         \quad(r=q,d).
 \end{aligned}\tag{RMW.5}
\]
The number of newly absent rational terms relative to (CAW.4) is
\[
                |Y_u|=N_T+J_+(T)+A_q(T)+A_d(T).
 \tag{RMW.6}
\]
The original triples no longer represented by a term of (RMW.3)
form the disjoint union $Z_u\sqcup V_u\sqcup Y_u$, of cardinality
$|Z_u|+2\mathbf1_{h\geq3}N_T+|Y_u|$, with $|Z_u|$ exactly
(Q7W.3). This counts represented triples, not zero positive indices.
For $T\geq2$ the indexing set of the rational sum is properly reduced;
no strict numerical decrease is asserted, since (RMW.4) is an equality.

Retain the exact costs $\mathcal E_u,\mathcal L_u,\mathcal A_u,
\mathcal P_u$ of (CAW.5), including $D_u=144\cdot2^h$ and the
distinct abelian degree $d_u=96\cdot2^h$.
With the same absolute constants as AC-G192-window,
\[
 \begin{aligned}
 |I_{399}-\mathcal R_{399}^{\rm match}|&\leq
       \mathcal E_u+\mathcal L_u+\mathcal A_u+C_2\mathcal P_u,\\
 |I_{399}^{\rm all7,2}-\mathcal R_{399}^{\rm match}|&\leq
       \mathcal E_u+\mathcal L_u+C_2\mathcal P_u.
 \end{aligned}\tag{RMW.7}
\]
There is one full (FL.4) local majorant, one (LT.2) torsion sum,
and one aggregate critical bound; the new families introduce
no additional cost or cardinality multiplier.
The uniform envelope remains
\[
 O_\varepsilon\bigl(u^{806+16\varepsilon}(\log u)^2\bigr)
                   =o\bigl(\exp(u^\varepsilon/2)\bigr).
 \tag{RMW.8}
\]
The exact positive-slack numerical premise for (RMW.3) supplies
the original full premise on the identical common window.
The preceding $\mathcal R_{399}^{\rm nc}$ premise is equivalent
with the identical threshold, window and slack by (RMW.4).
All remaining mixed and noncritical terms, and the retained
$\mathcal R_{\rm cub}$ and H04.4 inputs, stay unchanged.
\end{theorem}

\begin{proof}
\proofstep{1. Actual coefficients and literal full-local application.}
For $Y_+$ the finite character is trivial. At coefficient seven,
$\tau_7^e$ has order dividing six, already allowed by the original
character order $3\cdot2^h$. Its values are the specified roots
of unity in the unchanged $O_7$, not values chosen in another
completion. The host contains $\mathbb Q(\zeta_7)$.
For $r=q,d$, (RQM.10)--(RQM.11) identify $\lambda_r$ with
$\chi_{-r}$, of primitive conductor $r$, because $r\equiv3\pmod4$.
Its values are $\pm1$ in the original $O_r$, and its field is
exactly the actual auxiliary factor $\mathbb Q(\sqrt{-r})$.
No conductor-$4r$ character is substituted.

For every displayed triple, the rational finite character $\psi$
satisfies $\overline\psi=\overline c_\ell^{\,e}$ and is unramified
outside $\ell$. Thus $\psi c_\ell^{-e}$ is near one on this same
original ring. At odd $\ell$ its sign is even, so $\psi$ is
precisely the noncritical extension used in (CAW.4).
Theorem (RQM.2)--(RQM.5) gives
\[
             \mathcal S^{\mathbb Q}_{\ell,\psi_{\eta,e}^{\rm nc}}(e)=0
                     \qquad((\ell,\eta,e)\in Y_u).
 \tag{RMW.9}
\]
The coefficient $D_O=E_\ell/O_\ell$ is canonically the original
$O_\ell\otimes_{\mathbb Z_\ell}\mathbb Q_\ell/\mathbb Z_\ell$.
G195 Step4 proves the full strict comparison, including both
completion factors at split primes, all original inertia kernels,
and inflation through the same localization set. Its rank-two
induced lattice splits at these odd primes since two is a unit.
The original normalization $b_\ell$, ramification and codifferent
are unchanged. The identification with the critical partner
(RQM.5) does not make that partner zero.

\proofstep{2. Exhaustive overlap check with the preceding omitted set.}
All triples of $Y_u$ have odd coefficient different from three.
They also have coefficient different from $269$: that prime is
inert modulo three and is not in $Y_+$ or $Y_7$; it cannot equal
$q$ or $d$, since $269\equiv1\pmod4$.
They are therefore disjoint from every old coefficient-two,
coefficient-three and coefficient-$269$ row and from $V_u$.

At every other coefficient prime the complete preceding union
has exactly two possible kinds of row, by (OPW.3), (I2W.3) and
the later additions confined to coefficients two and three:
first, high-conductor two-power characters $\delta$ times a
modulo-seven character, with all auxiliary factors trivial;
second, the trivial label at an inert prime and a weight divisible
by $\ell-1$. In the first kind $\delta$ is nontrivial: its
primitive conductor exponent satisfies
$2^{c_\ell}>21\ell\,2^{2v_2(\ell^2-1)}$, hence $c_\ell>1$.
The labels in $Y_+$ and $Y_7$ have trivial two-power part.
Those in $Y_q,Y_d$ have a nontrivial auxiliary factor, absent
from every such old high-conductor row. None can meet that kind.
The second old kind cannot meet $Y_+$ because split and inert
primes are disjoint; it cannot meet $Y_7$ because seven is split;
and it cannot meet $Y_q,Y_d$ because their labels are nontrivial.
These alternatives exhaust $Z_u$ outside the three special
coefficient primes. Thus $Y_u\cap Z_u=\varnothing$.

No hidden character equality is introduced by restriction to $K$.
All the rational labels used in this overlap comparison factor
through $F_0$, whose intersection with $K$ is $\mathbb Q$.
Equivalently, their ratios have conductor prime to three and
cannot equal the nontrivial kernel character $\chi_{-3}$.
Coprime two, seven, $q$ and $d$ conductors distinguish the factors.
Every new label has trivial Kummer cubic factor. This proves
$Y_u\subset J_{u,0}\setminus V_u$, including the exact subtraction
of all possible old overlaps rather than an assumed new count.

\proofstep{3. Counts and equality of the explicit remainders.}
At a fixed split coefficient the positive multiples of $\ell-1$
below $T$ number $\lfloor(T-1)/(\ell-1)\rfloor$; for $\ell>T$
there are none. At coefficient seven there is exactly one label
at every weight. The trivial labels at weights divisible by six
are included there, and excluding seven from $\mathfrak P_+$
prevents their double counting. For $r=q,d$, put $a=(r-1)/2$.
The eligible weights are $a,a+2a,a+4a,\ldots$, giving
$\lfloor(T-1+a)/(2a)\rfloor=A_r(T)$, including zero when $T\leq a$.
The auxiliary coefficient primes are distinct, and their labels
are nontrivial even when those coefficients are split.
This proves the internal disjointness and (RMW.6).
For $T\geq2$ the actual triple $(7,\tau_7|_{G_K},1)$ is present,
so at least one previously retained rational term is removed.

All groups and sums are finite by AC-G192 and AC-G191.
By (RMW.9) each newly removed normalized logarithm equals zero.
Deleting exactly this finite disjoint family from (CAW.4) therefore
proves (RMW.4), with the original normalizations and no error.
Adding the three disjoint cardinalities gives the represented-triple
count in the theorem; the types of term being absent remain distinct.

\proofstep{4. One existing error, the same window and source rank.}
Substitute the exact equality (RMW.4) into the two inequalities
(CAW.6). This proves (RMW.7) without adding (RQM.6) or (CP.24)
a second time. All critical partners of $Y_u$ were already terms
of the same aggregate over $J_{u,0}$ in G192 Step3.
The one full local and one torsion majorant in G192 Step4 are
unchanged; no sum of separate family-majorants is used.
The retained estimate (CAW.7) proves (RMW.8).

Fix the original finite datum, page and $0<\varepsilon<1/256$.
Assume the (TC.3) premise with (RMW.3): there are $J\geq1$,
$X_0$ and $s\in(0,1)$, independent of $X,q,d,u,\ell,\eta,e$,
such that each prescribed pair at $X\geq X_0$ supplies one $W$
with $X^{1/J}\leq W$, $W^2\leq X$, and the stated bound
$\mathcal R_{399}^{\rm match}(u)\leq(1-s)\exp(u^\varepsilon/2)$
for every $u\in[W,W^2]$.
Choose a uniform $U_0$ making the first error in (RMW.7) at most
$(s/2)\exp(u^\varepsilon/2)$ for every permitted $u\geq U_0$.
At $X\geq\max(X_0,U_0^J)$ keep that exact $W$; the full original
numerical premise follows with slack $s/2$.
Between the old and new explicit remainders even this slack loss
is unnecessary, because they are equal pointwise.
The original actual rank $R.t$, simultaneous pages, coefficients,
local conditions and retained mixed realization are unchanged.
For a fixed finite simultaneous family take one maximum of
thresholds; no separate existential windows are intersected.
\end{proof}

\paragraph{Finite-cover dictionary and the remaining arithmetic gate.}
In G195's finite near-one argument, an original quotient
$D_O[\pi^m]$ with its continuous action unramified outside $\ell$
is a finite locally constant \emph{\'etale} coefficient module
on $\operatorname{Spec}\mathbb Z[1/\ell]$.
Its finite semidirect cocycle gives the corresponding finite
\'etale torsor; the local witness is retained in that same quotient.
The strict group is the kernel of the completion restriction at
$\ell$ and the original inertia restrictions elsewhere.
At the odd primes used here the real cohomology condition is zero.
This is the retained finite-cover/fundamental-group dictionary
of SGA1, Expos\'e V, \S\S4,7, applied to the exact original modules.
It introduces no new motive, polarization or integral comparison
with an analytic regulator. No fresh proof of SGA1 is claimed.

The remaining sum is still explicitly (RMW.3), with every mixed
cubic term and every unsupported noncritical rational term.
Its required uniform numerical bound, the mixed theta/Selmer
realization and H04.4 remain unproved. The new result removes
actual zero noncritical terms at growing coefficient primes and
at the actual auxiliary factors; it does not prove original-$K$
strict zero, individual positive-index zero, or strict numerical
decrease. G194 remains unused as already recorded in G192.
